\section{Consideraciones éticas}

El riesgo principal de Doctor Maíz no es que una métrica baje unas décimas,
sino que una persona convierta una predicción en una acción sobre el cultivo.
Un falso negativo puede retrasar la consulta cuando la planta presenta necrosis
letal o tizón. Un falso positivo de deficiencia puede inducir una aplicación de
fertilizante que no era necesaria. También puede existir una predicción
visualmente razonable pero agronómicamente incompleta: síntomas parecidos
pueden tener causas distintas y una fotografía no contiene información sobre
suelo, clima, etapa fenológica o evolución temporal.

\subsection{Confianza y comunicación del resultado}

La interfaz no presenta el resultado como un diagnóstico confirmado. Muestra
probabilidad, segunda clase cuando la confianza es menor de 0.70 y un mensaje
específico cuando el detector fuera de dominio rechaza la imagen. La demo
experimental emplea un umbral de 0.60 para advertir baja confianza y acompaña
cada salida con la frase «Apoyo orientativo; no sustituye la evaluación de un
profesional agrónomo». Estos umbrales son controles de interfaz, no garantías
estadísticas. No se ejecutó una calibración formal sobre el holdout y, por
tanto, 80\% de confianza no debe leerse como 80\% de certeza clínica o
agronómica.

El rechazo fuera de dominio de la aplicación usa distancia relativa de
Mahalanobis sobre un segundo tensor de características. Además, rechaza
distribuciones con poca separación. Esto reduce el riesgo de clasificar una
persona, una herramienta o una hoja de otra especie como una de las nueve
clases. No cubre todas las imágenes extrañas posibles y pertenece al modelo
móvil histórico, no al ensemble experimental. La selección de Etapa 2 todavía
no incorpora su propio detector OOD.

\subsection{Errores con costos distintos}

Las nueve clases no tienen el mismo costo de error. Confundir nitrógeno con
potasio puede terminar en una recomendación de manejo distinta; confundir una
hoja sana con una enfermedad puede causar gasto y preocupación; declarar sana
una hoja enferma puede demorar atención. El entrenamiento optimiza macro-F1,
que protege estadísticamente a las clases pequeñas, pero no incorpora una matriz
de costos agronómicos. Esa matriz debe definirse con especialistas y con
información de manejo local antes de convertir el modelo en un sistema de
decisión.

Por esa razón las recomendaciones de la app deben entenderse como pasos de
inspección: revisar otras hojas, observar la distribución del síntoma, repetir
la captura y consultar asistencia técnica. No corresponde indicar dosis o
productos solamente a partir del logit más alto. Una recomendación automática
de insumos sería una función nueva, con riesgos y validación distintos a los de
clasificar una imagen.

\subsection{Datos, ubicación y trazabilidad}

La base móvil almacena localmente ruta de imagen, etiqueta, confianza,
distribución y fecha. También reserva campos opcionales de latitud y longitud.
La operación offline evita enviar cada fotografía por defecto, pero no elimina
el riesgo de privacidad si más adelante se sincroniza una contribución. La
ubicación de una parcela puede revelar domicilio, patrón productivo o actividad
económica. Debe solicitarse consentimiento separado, explicar el propósito,
reducir precisión cuando no sea necesaria y permitir borrar el dato.

Cada predicción que se use para seguimiento debería conservar versión del
modelo, hash del artefacto, fecha y origen del resultado. La aplicación
inspeccionada valida el orden de etiquetas y el contrato de tensores; el
manifiesto del modelo registra arquitectura y SHA-256. Esa trazabilidad evita
un problema concreto: que una corrección del usuario quede asociada a un modelo
distinto del que produjo el diagnóstico.

\subsection{Supervisión y salida segura}

El sistema debe permitir «no sé». Una foto desenfocada, una hoja pequeña dentro
del encuadre o una condición que no pertenece a las nueve clases necesita
terminar en repetición de captura o consulta, no en una recomendación forzada.
El protocolo de campo propuesto incluye revisar varios órganos de la planta y
no actuar sobre una sola imagen. Las correcciones guardadas por productores o
extensionistas tampoco deben incorporarse automáticamente al entrenamiento:
primero requieren revisión agronómica, control de duplicados y registro de
procedencia.

Estas medidas no convierten el sistema en inocuo. Hacen visible dónde termina
la clasificación y dónde empieza una decisión humana. Ese límite es
especialmente necesario porque las métricas finales provienen de colecciones
públicas y no de una validación prospectiva en parcelas salvadoreñas.
